\subsection{Environmental Dependence and Low-Density Regimes}
  \label{subsec:environment}

  The effective modification introduced above is intrinsically sensitive to the
  gravitational environment.
  Regions of low baryonic density or weak gravitational gradients
  probe the saturated regime most strongly.

  This environmental dependence becomes critical at cosmological scales.
  Cosmic voids, characterized by extreme matter dilution,
  naturally sample the deep saturation regime.
  As a result, local kinematic measurements performed in void-dominated regions
  may exhibit systematic offsets relative to globally inferred expansion rates.

  Within this framework, the Hubble tension is not interpreted
  as a failure of early-universe cosmology.
  It arises as a late-time inference bias,
  reflecting the use of a single spacetime-based parametrization
  across environments characterized by different effective resolutions.

  Galaxy rotation curves and the Hubble tension therefore emerge
  as two manifestations of the same low-density phenomenology.
  Galaxies probe saturation radially through declining surface density,
  while cosmic voids probe it volumetrically through large-scale dilution.
  The identification of a single scale $a_\star$
  accounting for both effects constitutes the central organizing principle of this effective description.

  The following sections test this effective framework
  against galactic rotation curve data
  and explore its implications for local measurements
  of the Hubble constant.
